% This file is generated from the corresponding Obsidian note.
% Source: Teoricos/GS - Teo8.md
\chapter{Particiones de la unidad}
\label{ch:teo-8}
\section{Particiones de la Unidad}
\begin{bookremark}{Motivación}
La idea central es \textbf{pegar datos locales} definidos en abiertos y producir a partir de ellos objetos globales sobre toda la variedad.

\end{bookremark}
\begin{booktheorem}{}
La funcion $f:\mathbb{R}\rightarrow\mathbb{R}$ definida por

\[
f(x) = \begin{cases}0 & \text{if } x < 0 \\ e^{-1/x}  & \text{if } x \geq 0 \end{cases}
\]
es de clase $C^{\infty}$ y su serie de Taylor $s$ alrededor de $a=0$ es identicamente nula. En particular $s(x)\neq f(x)$ para todo $x>0$

\end{booktheorem}
\phantomsection\label{blk:54d8a6}
\section{Funciones de corte y funciones campana}
\begin{booklemma}{Función de corte en $\mathbb R$}
Dados $\gamma_1,\gamma_2\in\mathbb R$ con $\gamma_1<\gamma_2$, existe una función suave

\[
h:\mathbb R\to\mathbb R
\]
tal que

\[
h(t)=\begin{cases} 1 & \text{if } t \le \gamma_1,\\ 0<h(t)<1 & \text{if } \gamma_1 < t < \gamma_2,\\ 0 & \text{if } t \ge \gamma_2 \end{cases}
\]

\bookimage{assets/pasted-image-20260414194706.png}{Pasted image 20260414194706}


\end{booklemma}
\begin{bookproof}{}
\begin{enumerate}
\item Considerar \hyperref[blk:54d8a6]{Teórico 8} y definir
\end{enumerate}
\[
h(t)= \frac{f(\gamma_{2}-t)}{f(\gamma_{2}-t)+f(t-\gamma_{1})}
\]
\bookqed
\end{bookproof}
\phantomsection\label{blk:3df6bd}
\begin{booklemma}{Función campana en $\mathbb R^n$}
Dados $0<\gamma_1<\gamma_2$, existe una función suave

\[
\mu:\mathbb R^n\to\mathbb R
\]
tal que

\[
\mu(x)= \begin{cases}  1 & \text{if } x\in \overline{B(0,\gamma_1)},\\ 0<\mu(x)<1 & \text{if } \gamma_1<\|x\|<\gamma_2,\\ 0 & \text{if } x\in B(0,\gamma_2)^c \end{cases}
\]

\bookimage{assets/pasted-image-20260414195214.png}{Pasted image 20260414195214}


\end{booklemma}
\begin{bookproof}{}
2. Se toma la función del lema anterior y se compone con la norma:

\[
\mu(x)=h(\|x\|).
\]
\begin{enumerate}[start=3]
\item La única sutileza es el punto $x=0$, pero cerca de $0$ la función es constante igual a $1$, así que es suave.
\end{enumerate}
\bookqed
\end{bookproof}
\phantomsection\label{blk:83d0e9}
\begin{bookexercise}{}
Sea $X$ un espacio topológico Hausdorff, $U\subset X$ con la topología subespacio, y $A\subset U$. Si existe un conjunto compacto $K$ tal que $A\subset K\subset U$ entonces

\[
\text{Clausura}_{U} A=\text{Clausura}_{X} A
\]
No olvidar que en general $\operatorname{Clausura}_{U}A \subseteq \operatorname{Clausura}_{X}A$ pero no tienen por que ser iguales ejemplo $U=A=(0,1)$ y $X=\mathbb{R}$

\end{bookexercise}
\begin{bookproof}{}
\begin{enumerate}
\item Como $K\subset U\subset X$ y $K$ es compacto en $U$, entonces también es compacto en $X$.
\item Como $X$ es Hausdorff, los compactos son cerrados. Luego, $K$ es cerrado en $X$
\item Dado que $A\subset K$ y $K$ es cerrado en $X$, se tiene $\operatorname{Cl}_{U} A\subset K\subseteq U$.
\item Por la definición de topología subespacio, $\operatorname{Cl}_{U} A=U\cap\operatorname{Cl}_{X} A$.
\item Como $K$ cerrado en $X$ entonces $\operatorname{Cl}_{X} A\subset K\subset U$, resulta $U\cap\operatorname{Cl}_{X} A=\operatorname{Cl}_{X} A$.
\end{enumerate}
\bookqed
\end{bookproof}
\section{No se por que aparecen devuelta si ya estaban antes}
\begin{booklemma}{Funciones campana o de levantamiento en variedades suaves}
Sea $M$ una variedad suave, $p\in M$ y $U$ un abierto de $M$ con $p\in U$. Entonces existen un abierto $V$ de $p$ y una función suave

\[
\beta:M\to\mathbb R
\]
tales que

\[
\overline V\subseteq U,\qquad 0\le \beta\le 1,\qquad \beta|_{\overline V}\equiv 1,\qquad \operatorname{supp}\beta\subseteq U.
\]
notar que aca la clausura es con respecto a $M$

\end{booklemma}
\begin{bookproof}{}
6. Sea $p\in M$ y $U$ abierto de $p$. Sea $(W,\psi)$ una carta suave de $p$ tal que, sin pérdida de generalidad, podemos asumir que:

\begin{itemize}
\item $W\subset U$
\item $\psi(W)=B(0,3)$ (Esto por que las bolas son base $\mathbb{R}^{n}$ y son difeomorfas entre si)
\end{itemize}
\begin{enumerate}
\item Por tanto hay un abierto $V$ de $p$, $V\subset W\subset U$ tal que:
\end{enumerate}
\[
\psi(V)=B(0,1)
\]
\begin{enumerate}[start=2]
\item Pensamos en la función $\beta:M\to\mathbb{R}$:
\end{enumerate}
\[
\beta(q)=\begin{cases}h\circ\psi(q) & \text{si } q\in W,\\0 & \text{si } q\notin W.\end{cases}
\]
con $h$ como \hyperref[blk:83d0e9]{Teórico 8} con $\gamma_1=1$, $\gamma_2=2$.\\
3. Veamos que cumple las condiciones que buscamos:

\begin{itemize}
\item (i)  $0\le \beta(x)\le 1$ es trivial
\item (ii)

\begin{enumerate}
\item Veamos $\beta|_{\overline{V}}\equiv 1$. Tenemos que $V=\psi^{-1}(B(0,1))$ que es un abierto de $W$ y por tanto es abierto de $M$.
\item Notemos que como $V\subseteq K=\psi ^{-1}(\overline{B(0,1)})\subseteq W$ con $K$ compacto por que $\psi$ es homeomorfismo entonces $\overline{V}^{W}=\overline{V}^{M}=\overline{V}$
\item Ahora como $\psi$ es homeo
\end{enumerate}
\end{itemize}
\[
\overline{V} =\overline{V}^{W}=\psi ^{-1}(\psi(\overline{V}^{B(0,3)}  ))=\psi ^{-1}(\overline{B(0,1)}^{B(0,3)})  \subseteq B(0,2)\subseteq W
\]
\begin{Verbatim}[fontsize=\small]
 4. En particular, $\overline{V}\subseteq W\subseteq U$ entonces $\beta|_{\overline{V}}\equiv 1$.  
\end{Verbatim}
\begin{itemize}
\item (ii) Soporte de $\beta\subseteq U$:

\begin{enumerate}
\item Recordemos que $\operatorname{supp}(\beta)=\overline{\{q\in M:\beta(q)\neq 0\}}$.
\item En el caso de la función $h$, el soporte es $\overline{B(0,2)}$.
\item Afirmación: $\operatorname{supp}(\beta)=\psi^{-1}(\overline{B(0,2)})$.
\item En efecto, por definición de $\beta$, $\beta(q)\neq 0$ sii $q\in W$ y $\psi(q)\in B(0,2)$.
\item Luego por definicion
\end{enumerate}
\end{itemize}
\[
\operatorname{supp}(\beta)=\overline{(\psi^{-1}(B(0,2)))}^{M}
\]
\begin{Verbatim}[fontsize=\small]
 6. Como $\psi^{-1}(B(0,2))$ está contenido en el compacto $K=\psi^{-1}(\overline{B(0,2.5)})$, que a su vez está contenido en $W$ devuelta las clausura coinciden 
\end{Verbatim}
\[
\overline{(\psi^{-1}(B(0,2)))}^{M}=\overline{(\psi^{-1}(B(0,2)))}^{W}
\]
\begin{Verbatim}[fontsize=\small]
 7. Por lo tanto se tiene: 
\end{Verbatim}
\[
\operatorname{supp}(\beta)=\overline{\psi ^{-1}(B(0,2)) }^{W}=\psi^{-1}(\overline{B(0,2)}^{B(0,3)}) \subset K\subset W\subset U
\]
\begin{Verbatim}[fontsize=\small]
 8. Nuevamente, hay que tener cuidado al clausurar: podría haber pasado algo como $U=(0,5)$ y $\psi^{-1}(\Delta(0,2))=(0,5)$, y en este caso se sale de $U$. Por eso tomamos $W$ tal que $\psi(W)=\Delta(0,3)$, para quedar holgados y que la clausura del soporte no se escape de $U$.  
 9. Por último, veamos que $\beta$ es suave mostrando que lo es en algún entorno abierto de cada punto.  
 10. Tenemos dos casos:  
 	- (i) Si $q\in W$, entonces $\beta|_{W}=h\circ\psi$, que es composición de suaves.
 	- (ii) Si $q\notin W$, entonces $q\not\in \operatorname{supp}(\beta)$ por que $\operatorname{supp}(\beta) \subseteq W$ y como $\operatorname{supp}(\beta)$ es cerrado, su complemento es abierto entonces existe un abierto $O\subseteq \operatorname{supp}(\beta)^{c}$ de $q$ tal que $\beta|_{O}=0$, como función constante es suave, se tiene que $\beta|_{O}$ es suave. 
 6. Por lo tanto, $\beta$ es suave.
\end{Verbatim}
\bookqed
\end{bookproof}
\begin{booklemma}{Extensión local de funciones suaves}
Sea $M$ una variedad suave, $p\in M$ y $U$ un abierto de $p$. Para toda

\[
f\in C^\infty(U)
\]
existen un abierto $V$ de $p$ con $\overline V\subseteq U$ y una función

\[
\widetilde f\in C^\infty(M)
\]
tal que

\[
\widetilde f=f\quad\text{sobre }V
\]
\end{booklemma}
\begin{bookproof}{}
7. Tomamos una función campana $\beta$ como en el lema anterior, con $\beta\equiv 1$ sobre $\overline V$ y $\operatorname{supp}\beta\subseteq U$. 8. Definimos

\[
\widetilde f(q)=\begin{cases}\beta(q)f(q),& q\in U,\\0,& q\notin U\end{cases}
\]
\begin{enumerate}[start=9]
\item Entonces $\widetilde f$ es suave, coincide con $f$ sobre $V$ y está bien definida globalmente.
\end{enumerate}
\bookqed
\end{bookproof}
\section{Familias localmente finitas}
\begin{bookdefinition}{Familia localmente finita}
Sea $X$ un espacio topológico. Una familia $\Omega=\{A_\alpha\}_{\alpha\in I}$ de subconjuntos de $X$ se dice \textbf{localmente finita} si para todo $p\in X$ existe un abierto $W$ de $p$ que corta a lo sumo una cantidad finita de miembros de $\Omega$.

\end{bookdefinition}
\begin{bookexample}{Ejemplos}
\begin{itemize}
\item La familia
\end{itemize}
\[
\bigl\{(k,k+2):k\in\mathbb N\bigr\}
\]
de abiertos de $\mathbb R$ es localmente finita.

\begin{itemize}
\item La familia
\end{itemize}
\[
\bigl\{(-1/k,1/k):k\in\mathbb N\bigr\}
\]
no es localmente finita, porque todo entorno de $0$ corta infinitos términos.

\end{bookexample}
\begin{bookdefinition}{Cubrimiento y refinamiento}
Sea $\Omega=\{C_\alpha\}_{\alpha\in I}$ una familia de subconjuntos de un espacio topológico $X$.

\begin{itemize}
\item $\Omega$ \textbf{cubre} a $X$ si
\end{itemize}
\[
\bigcup_{\alpha\in I} C_\alpha=X.
\]
\begin{itemize}
\item Un \textbf{refinamiento} de $\Omega$ es otro cubrimiento
\end{itemize}
\[
\mathcal D=\{D_\beta\}_{\beta\in J}
\]
tal que para todo $\beta\in J$ existe $\alpha\in I$ con

\[
D_\beta\subseteq C_\alpha.
\]
\end{bookdefinition}
\section{Particiones de la unidad}
\begin{bookdefinition}{Partición de la unidad}
Una \textbf{partición de la unidad} sobre una variedad suave $M$ es una familia de funciones suaves

\[
\{\rho_\alpha:M\to\mathbb R\}_{\alpha\in I}
\]
tal que:

\begin{enumerate}
\item la familia de soportes $\{\operatorname{supp}\rho_\alpha\}_{\alpha\in I}$ es localmente finita;
\item para todo $\alpha$ y todo $p\in M$,
\end{enumerate}
\[
0\le \rho_\alpha(p)\le 1;
\]
\begin{enumerate}[start=3]
\item para todo $p\in M$,
\end{enumerate}
\[
\sum_{\alpha\in I}\rho_\alpha(p)=1.
\]
La suma está bien definida porque la familia de soportes es localmente finita.

\end{bookdefinition}
\phantomsection\label{blk:f5db1b}
\begin{bookdefinition}{Subordinación}
Sea $\Omega=\{U_\alpha\}_{\alpha\in I}$ un cubrimiento por abiertos de $M$.

\begin{itemize}
\item Se dice que la particion de la unidad $\{\rho_j\}_{j\in J}$ esta subordinada a $\Omega$ si la familia de soportes $\{ \operatorname{supp}(\rho_j) \}_{j\in J}$ es un refinamiento de $\Omega$.
\item Y se dice que la partición de la unidad $\{\rho_j\}_{j\in J}$ esta \textbf{estrictamente subordinada} a $\Omega$ si $J=I$ y $\operatorname{supp}\rho_\alpha\subseteq U_\alpha$ para todo $\alpha\in I$. Notar que siempre una particion de la unidad cubre todo $M$ por parte 3 de la dfeinicion.
\end{itemize}
\end{bookdefinition}
\phantomsection\label{blk:12d7c6}
\section{Ejemplo (Aplicaciones)}
\begin{bookremark}{}
Particiones de la unidad sirven para pegar funciones suaves definidas en abiertos pequeños de la variedad. Supongamos que tenemos una familia de funciones suaves $\{f_\alpha:U_\alpha\subseteq M\to\mathbb{R}\}_{\alpha\in I}$ con $\{U_\alpha\}_{\alpha\in I}$ un cubrimiento abierto de $M$, y que además tenemos una partición de la unidad $\{\rho_\alpha\}_{\alpha\in I}$ estrictamente subordinada a $\{U_\alpha\}_{\alpha\in I}$. Entonces,

\[
f:=\sum_{\alpha\in I}\rho_\alpha f_\alpha
\]
es una función suave en $C^\infty(M)$, donde el producto $\rho_\alpha\cdot f_\alpha$ se ve como la función suave definida sobre TODO $M$ $(\leftarrow$ abuso de notación$)$ dada por

\[
(\rho_\alpha\cdot f_\alpha)(q)=\begin{cases}\rho_\alpha(q)\,f_\alpha(q)&\text{si }q\in U_\alpha,\\0&\text{si }q\notin U_\alpha.\end{cases}
\]
que resulta suave pues $\operatorname{supp}(\rho_\alpha)\subseteq U_\alpha$. Y como la familia de soportes es localmente finita, para cada $p\in M$, existe un abierto de $p$ en donde la suma anterior se vuelve una suma finita de funciones suaves.

\end{bookremark}
\begin{bookremark}{Posible para final}
Si $M$ es compacta, es fácil dar particiones de la unidad subordinadas a cualquier cubrimiento abierto $\{U_\alpha\}_{\alpha\in I}$ de $M$. Por \hyperref[blk:45cd5d]{Teórico 5}, dado $p\in U_\alpha$, existe un abierto $V_p$ de $p$ t.q. $\overline{V_p}\subseteq U_\alpha$ y una función $B_p:M\to\mathbb{R}$ suave t.q.

\[
\begin{cases}0\leq B_p(q)\leq 1,&\forall q\in M,\\B_p(q)=1,&\forall q\in \overline{V_p},\\\operatorname{supp}B_p\subseteq U_\alpha.\end{cases}
\]
Dado que $\{V_p\}_{p\in M}$ cubre a $M$, se sigue por compacidad de $M$ que existen puntos $p_1,\dots,p_r$ t.q. $\{V_{p_k}:k=1,\dots,r\}$ cubre a $M$. Ahora definir

\[
\rho_k=\frac{B_{p_k}}{\sum_{k=1}^r B_{p_k}},\qquad k=1,\dots,r
\]
que resulta suave pues el denominador nunca se anula y $\{ \rho_{k} \}_{k=1}^{n}$ es una particion de la unidad subordinada a $\{ U_{\alpha } \}_{\alpha \in J}$

\end{bookremark}
\begin{bookremark}{}
La construccion anterior se generaliza a lo que se conoce como \textbf{espacios paracompactos}

\end{bookremark}
\section{Espacios paracompactos}
\begin{bookdefinition}{Espacio paracompacto}
Un espacio topológico $X$ se dice \textbf{paracompacto} si todo cubrimiento abierto de $X$ admite un refinamiento abierto que es localmente finito.

\end{bookdefinition}
\begin{booktheorem}{Las variedades topológicas son paracompactas}
Toda variedad topológica es paracompacta. En particular, toda variedad suave es paracompacta.

\end{booktheorem}
\begin{booktheorem}{Existencia de particiones de la unidad}
Sea $M$ una variedad suave y sea

\[
\Omega=\{U_\alpha\}_{\alpha\in I}
\]
un cubrimiento abierto de $M$. Entonces existe una partición de la unidad \textbf{numerable}

\[
\{\rho_k\}_{k\in\mathbb N}
\]
subordinada a $\Omega$ tal que cada $\operatorname{supp}\rho_k$ es compacto para todo $k\in \mathbb{N}$. Más aún, si no se exige compacidad de los soportes, puede construirse una partición estrictamente subordinada al $\Omega$ donde a lo sumo una cantidad numerable de los $\rho_{\alpha }$ son no nulos

\end{booktheorem}
\phantomsection\label{blk:eafe2b}
\begin{bookexample}{Caso compacto}
Si $M$ es compacta, toda cobertura abierta admite un subcubrimiento finito

\[
U_1,\dots,U_r.
\]
Tomando funciones campana $\beta_i$ con soporte contenido en $U_i$ y tales que los abiertos donde $\beta_i=1$ aún cubren $M$, se define

\[
\rho_i=\frac{\beta_i}{\beta_1+\cdots+\beta_r}.
\]
Entonces $\{\rho_i\}_{i=1}^r$ es una partición de la unidad subordinada al cubrimiento.

\end{bookexample}
\section{Consecuencias de las particiones de la unidad}
\begin{bookcorollary}{Función que vale 1 sobre un cerrado}
Sea $U$ un abierto de $M$ y $A$ un cerrado de $M$ con

\[
A\subseteq U
\]
Entonces existe una función suave

\[
f:M\to\mathbb R
\]
tal que

\[
0\le f\le 1,\qquad f|_A\equiv 1,\qquad \operatorname{supp}f\subseteq U
\]
\end{bookcorollary}
\begin{bookproof}{}
\begin{enumerate}
\item Como $A$ es cerrado. Considerar al cubrimiento de $M$ dado por
\end{enumerate}
\[
\{U,M\setminus A\}
\]
\begin{enumerate}[start=2]
\item Por \hyperref[blk:eafe2b]{Existencia de particiones de la unidad} existe una particion de la unidad $\{\rho,\sigma\}$ subordinada a $\{U,M\setminus A\}$ con
\end{enumerate}
\[
\operatorname{supp}\rho\subseteq U,\qquad \operatorname{supp}\sigma\subseteq M\setminus A
\]
\begin{enumerate}[start=3]
\item Para $p\in A$ se tiene $\sigma(p)=0$ porque $\operatorname{supp}\sigma\subseteq M\setminus A$. Usando la parte 3 de la \hyperref[blk:f5db1b]{definición de partición de la unidad}, tenemos $\rho(p)+\sigma(p)=1$, luego
\end{enumerate}
\[
\rho(p)=1
\]
\begin{enumerate}[start=4]
\item Tomamos $f=\rho$ que es la funcion que buscabamos
\end{enumerate}
\bookqed
\end{bookproof}
\begin{bookdefinition}{Función suave sobre un subconjunto arbitrario}
Sean $M$ y $N$ variedades suaves y sea $A\subseteq M$ un subconjunto cualquiera. Una función

\[
F:A\to N
\]
se dice \textbf{suave} si para todo $p\in A$ existe un abierto $W_p$ de $M$ con $p\in W_p$ y una función suave

\[
F_p:W_p\to N
\]
que coincide con $F$ sobre $W_p\cap A$

\end{bookdefinition}
\phantomsection\label{blk:cf02ec}
\begin{bookproposition}{Extensión de funciones suaves definidas en cerrados}
Sea $M$ una variedad suave, $A\subseteq M$ cerrado y

\[
F:A\to\mathbb R^k
\]
una función suave en el sentido anterior. Entonces, para todo abierto $U$ con $A\subseteq U$ existe una función suave

\[
\widetilde F:M\to\mathbb R^k
\]
funcion suave tal que

\[
\widetilde F|_A=F,
\qquad
\operatorname{supp}\widetilde F\subseteq U.
\]
\end{bookproposition}
\begin{bookproof}{}
\begin{enumerate}
\item Por \hyperref[blk:cf02ec]{Función suave sobre un subconjunto arbitrario}, para cada $p\in A$ elegimos un abierto $W_p\subseteq U$ y una extensión suave local
\end{enumerate}
\[
\widetilde F_p:W_p\to\mathbb R^k
\]
que coincide con $F$ en $W_{p}\cap A$ 2. Sin perdida de generalidades podemos suponer $W_{p}\subseteq U$ (En otro caso, cambiamos $W_{p}$ por el abierto $W_{p}\cap U$) 3. Ahora tenemos el conjunto

\[
\{W_p\}_{p\in A}\cup\{M\setminus A\}
\]
que es un cubrimiento abierto de $M$. 4. Por \hyperref[blk:eafe2b]{Existencia de particiones de la unidad} existe una partición de la unidad

\[
\{\rho_p\}_{p\in A}\cup\{\rho_0\}
\]
subordinada estrictamente a este cubrimiento, con

\[
\operatorname{supp}\rho_p\subseteq W_p, \qquad \operatorname{supp}\rho_0\subseteq M\setminus A
\]
\begin{enumerate}[start=5]
\item Como hicimos en el ejemplo anterior, definimos
\end{enumerate}
\[
\widetilde F=\sum_{p\in A}\rho_p\,\widetilde F_p
\]
donde estamos abusando de notación y entendemos el producto

\[
(\rho_p\widetilde{F}_p)(q)=\begin{cases}\rho_p(q)\,\widetilde{F}_p(q)&\text{si }q\in W_p,\\0&\text{si }q\notin W_p.\end{cases}
\]
\begin{enumerate}[start=6]
\item Veamos que cada término $\rho_p\widetilde F_p$ extendido por $0$ fuera de $W_p$ es suave sobre todo $M$. Si $q\in W_p$, la suavidad es trivial porque allí es producto de funciones suaves. Si $q\notin W_p$, entonces como $\operatorname{supp}\rho_p\subseteq W_p$, también $q\notin\operatorname{supp}\rho_p$. Como $\operatorname{supp}\rho_p$ es cerrado, $M\setminus\operatorname{supp}\rho_p$ es un abierto que contiene a $q$, y en ese abierto la extensión vale $0$; por tanto también es suave alrededor de $q$.
\item Tenemos que $\widetilde{F}$ es suave: como $\{\operatorname{supp}\rho_p\}$ es localmente finita, alrededor de cada punto de $M$ la suma es finita; y por el paso 6 cada término de esa suma finita es una función suave.
\item Ahora, si $q\in A$,
\end{enumerate}
\[
\widetilde{F}(q)=\sum_{p\in A}\rho_p(q)\widetilde{F}_p(q)=\sum_{p\in A}\rho_p(q)F(q)
\]
porque $\widetilde{F}_p$ coincide con $F$ en $W_p\cap A$, 9. Y seguimos

\[
\sum_{p\in A}\rho_p(q)F(q)=\left(\sum_{p\in A}\rho_p(q)\right)F(q)=\left(1-\rho_0(q)\right)F(q)
\]
donde el segundo igual vale por que $\sum_{p\in A}\rho_{p}+\rho_{0}=1$ 10. Y como $\operatorname{supp}\rho_0\subseteq M\setminus A$, tenemos

\[
\left(1-\rho_0(q)\right)F(q)=F(q).
\]
\begin{enumerate}[start=11]
\item Por último veamos $\operatorname{supp}\widetilde{F}\subseteq U$: Sea $q\in\operatorname{supp}\widetilde{F}$. Como $\{\operatorname{supp}\rho_p\}\cup\{\operatorname{supp}\rho_0\}$ es localmente finita, existe $W_q$ abierto de $M$ con $q$ que corta un número finito de miembros de $\{\operatorname{supp}\rho_p\}\cup\{\operatorname{supp}\rho_0\}$, digamos $\operatorname{supp}\rho_{p_1},\ldots,\operatorname{supp}\rho_{p_k}$ y $\operatorname{supp}\rho_0$.
\item Esto implica que debe haber un $j$ tal que $q\in\operatorname{supp}\rho_{p_j}$, pues en otro caso, existirían abiertos $W_1,\ldots,W_k$ de $q$ donde $\rho_{p_1},\ldots,\rho_{p_k}$ se anulan, y así
\end{enumerate}
\[
W_0\cap W_1\cap\cdots\cap W_k
\]
es un abierto de $q$ donde $\widetilde{F}$ vale cero que es absurdo 13. Como $q\in\operatorname{supp}\rho_{p_j}$ y $\operatorname{supp}\rho_{p_j}\subseteq W_{p_j}\subseteq U$, entonces $q\in U$. Por lo tanto, $\operatorname{supp}\widetilde F\subseteq U$.

\bookqed
\end{bookproof}
\phantomsection\label{blk:b5794a}
\section{Métrica riemanniana}
\begin{bookdefinition}{Métrica riemanniana}
Sea $M$ una variedad suave. Una \textbf{métrica riemanniana} $g$ sobre $M$ es (por ahora) una asignación

\[
g:p\mapsto g_p,
\]
donde cada $g_p$ es un producto interno sobre $T_pM$, tal que si

\[
(U,\varphi=(x_1,\dots,x_n))
\]
es una carta suave de $M$, entonces las funciones

\[
g_{ij}:U\to\mathbb R,
\qquad
p\mapsto g_p\left(\left.\frac{\partial}{\partial x_i}\right|_p,\left.\frac{\partial}{\partial x_j}\right|_p\right)
\]
son suaves para todo $i,j$.

\end{bookdefinition}
\begin{bookexercise}{Estructuras riemannianas}
Mostrar que toda variedad diferenciable admite una estructura riemanniana. Más precisamente, probar que existe una aplicación

\[
B:p\mapsto B_p
\]
que asigna a cada $p\in M$ un producto interno en $T_pM$ y tal que sus coeficientes locales son suaves.

¿Cómo se definiría, a partir de una métrica riemanniana, la longitud de una curva diferenciable

\[
\alpha:[0,1]\to M?
\]
Pista: usar particiones de la unidad.

\end{bookexercise}
